\documentclass[11pt,a4paper]{article}
\usepackage[utf8]{inputenc}
\usepackage{amsmath,amssymb,amsthm}
\usepackage{algorithm}
\usepackage{algpseudocode}
\usepackage{booktabs}
\usepackage{hyperref}
\usepackage{geometry}
\geometry{margin=1in}

\title{Long Short-Term Memory (LSTM): Constant Error Carousel, Gated State Transitions, and Long-Horizon Gradient Propagation}
\author{Automorphic Intelligence Architecture Specifications}
\date{\today}

\newtheorem{theorem}{Theorem}
\newtheorem{lemma}[theorem]{Lemma}
\newtheorem{proposition}[theorem]{Proposition}
\theoremstyle{definition}
\newtheorem{definition}{Definition}
\newtheorem{remark}{Remark}

\begin{document}

\maketitle

\begin{abstract}
Standard recurrent neural architectures fail to bridge extended temporal delays due to exponential gradient decay during backpropagation. Long Short-Term Memory (LSTM) overcomes this limitation by introducing a dedicated internal memory cell state governed by an additive Constant Error Carousel (CEC) modulated by continuous multi-plicative gating mechanisms (forget, input, and output gates). This document presents the complete mathematical state-space formulation, proves gradient preservation through the additive cell recurrence path, details algorithmic pseudocode, and provides a verified step-by-step numerical calculation.
\end{abstract}

\section{Notation and Mathematical Preliminaries}

Let $\mathbf{x}_t \in \mathbb{R}^{d_x}$ denote the input vector at time $t$, $\mathbf{h}_t \in \mathbb{R}^{d_h}$ denote the external hidden state, and $\mathbf{c}_t \in \mathbb{R}^{d_h}$ denote the internal cell memory state.

\begin{table}[h]
\centering
\caption{Mathematical Notation for LSTM Architecture}
\begin{tabular}{lll}
\toprule
\textbf{Symbol} & \textbf{Domain} & \textbf{Semantic Description} \\
\midrule
$\mathbf{f}_t$ & $(0, 1)^{d_h}$ & Forget gate activation vector \\
$\mathbf{i}_t$ & $(0, 1)^{d_h}$ & Input gate activation vector \\
$\tilde{\mathbf{c}}_t$ & $(-1, 1)^{d_h}$ & Candidate memory update vector \\
$\mathbf{o}_t$ & $(0, 1)^{d_h}$ & Output gate activation vector \\
$W_{\{f,i,c,o\}}$ & $\mathbb{R}^{d_h \times d_x}$ & Input projection weight matrices \\
$U_{\{f,i,c,o\}}$ & $\mathbb{R}^{d_h \times d_h}$ & Recurrent hidden projection weight matrices \\
$\mathbf{b}_{\{f,i,c,o\}}$ & $\mathbb{R}^{d_h}$ & Gate bias vectors \\
$\sigma(\cdot)$ & $\mathbb{R} \to (0, 1)$ & Elementwise logistic sigmoid activation \\
$\odot$ & Operator & Hadamard elementwise product \\
\bottomrule
\end{tabular}
\end{table}

\section{Problem Definition and Core Concepts}

\subsection{3-Level Explanation}
\begin{enumerate}
    \item \textbf{Intuitive (High School level):} Unlike a regular RNN that scrambles its thoughts every step, an LSTM has a notebook (the cell). Three security guards control the notebook: the Forget Guard erases useless old notes, the Input Guard writes down new facts, and the Output Guard reads out what to say.
    \item \textbf{Engineering Level:} LSTM replaces multiplicative matrix-product recurrence with an additive conveyor belt $\mathbf{c}_t = \mathbf{f}_t \odot \mathbf{c}_{t-1} + \mathbf{i}_t \odot \tilde{\mathbf{c}}_t$. If the forget gate $\mathbf{f} \approx \mathbf{1}$, gradients flow across hundreds of time steps without decaying, enabling long-range dependency modeling.
    \item \textbf{Mathematical Level:} Given sequence $\{\mathbf{x}_t\}$, the state transition system is parameterized by:
    \begin{align}
    \mathbf{f}_t &= \sigma(W_f \mathbf{x}_t + U_f \mathbf{h}_{t-1} + \mathbf{b}_f) \\
    \mathbf{i}_t &= \sigma(W_i \mathbf{x}_t + U_i \mathbf{h}_{t-1} + \mathbf{b}_i) \\
    \tilde{\mathbf{c}}_t &= \tanh(W_c \mathbf{x}_t + U_c \mathbf{h}_{t-1} + \mathbf{b}_c) \\
    \mathbf{c}_t &= \mathbf{f}_t \odot \mathbf{c}_{t-1} + \mathbf{i}_t \odot \tilde{\mathbf{c}}_t \\
    \mathbf{o}_t &= \sigma(W_o \mathbf{x}_t + U_o \mathbf{h}_{t-1} + \mathbf{b}_o) \\
    \mathbf{h}_t &= \mathbf{o}_t \odot \tanh(\mathbf{c}_t)
    \end{align}
\end{enumerate}

\subsection{5-Step Algorithmic Framework}
\begin{enumerate}
    \item \textbf{Fused Gate Pre-activation:} Compute linear combinations $[W_f; W_i; W_c; W_o] \mathbf{x}_t + [U_f; U_i; U_c; U_o] \mathbf{h}_{t-1} + \mathbf{b}$ in a single matrix multiplication.
    \item \textbf{Non-linear Gating Activation:} Apply sigmoid $\sigma$ to gates $(\mathbf{f}_t, \mathbf{i}_t, \mathbf{o}_t)$ and hyperbolic tangent to candidate $\tilde{\mathbf{c}}_t$.
    \item \textbf{Additive Cell State Evolution:} Modulate prior cell $\mathbf{c}_{t-1}$ by forget gate and add gated candidate $\mathbf{i}_t \odot \tilde{\mathbf{c}}_t$.
    \item \textbf{Gated Hidden Emission:} Compute $\mathbf{h}_t = \mathbf{o}_t \odot \tanh(\mathbf{c}_t)$.
    \item \textbf{Temporal State Roll-Forward:} Pass $(\mathbf{h}_t, \mathbf{c}_t)$ to step $t+1$.
\end{enumerate}

\section{Algorithm Specification}

\begin{algorithm}[H]
\caption{Long Short-Term Memory Step and Unroll}
\label{alg:lstm}
\begin{algorithmic}[1]
\Require Sequence $\{\mathbf{x}_t\}_{t=1}^T$, initial states $\mathbf{h}_0 = \mathbf{0}, \mathbf{c}_0 = \mathbf{0}$, weights $\mathbf{W}, \mathbf{U}$, biases $\mathbf{b}$.
\Ensure Hidden state trajectory $\{\mathbf{h}_t\}_{t=1}^T$, cell state trajectory $\{\mathbf{c}_t\}_{t=1}^T$.
\For{$t = 1$ \textbf{to} $T$}
    \State $\mathbf{z} \gets \mathbf{W} \mathbf{x}_t + \mathbf{U} \mathbf{h}_{t-1} + \mathbf{b}$
    \State $\mathbf{f}_t \gets \sigma(\mathbf{z}[0 : d_h])$
    \State $\mathbf{i}_t \gets \sigma(\mathbf{z}[d_h : 2d_h])$
    \State $\tilde{\mathbf{c}}_t \gets \tanh(\mathbf{z}[2d_h : 3d_h])$
    \State $\mathbf{o}_t \gets \sigma(\mathbf{z}[3d_h : 4d_h])$
    \State $\mathbf{c}_t \gets \mathbf{f}_t \odot \mathbf{c}_{t-1} + \mathbf{i}_t \odot \tilde{\mathbf{c}}_t$
    \State $\mathbf{h}_t \gets \mathbf{o}_t \odot \tanh(\mathbf{c}_t)$
\EndFor
\State \Return $\{\mathbf{h}_t\}_{t=1}^T, \{\mathbf{c}_t\}_{t=1}^T$
\end{algorithmic}
\end{algorithm}

\section{Theoretical Guarantees and Constant Error Flow}

\begin{theorem}[Constant Error Carousel Gradient Preservation]
\label{thm:cec_gradient}
Let the forget gate be saturated to $\mathbf{f}_k = \mathbf{1}$ for all steps $k \in [t+1, T]$. The temporal gradient of the cell state $\frac{\partial \mathbf{c}_T}{\partial \mathbf{c}_t}$ satisfies:
\begin{equation}
\frac{\partial \mathbf{c}_T}{\partial \mathbf{c}_t} = \prod_{k=t+1}^T \text{diag}(\mathbf{f}_k) = I_{d_h}
\end{equation}
guaranteeing zero exponential vanishing of error signals along the cell trajectory over arbitrary temporal spans $(T - t)$.
\end{theorem}

\begin{proof}
Directly differentiating the linear cell state recurrence $\mathbf{c}_k = \mathbf{f}_k \odot \mathbf{c}_{k-1} + \mathbf{i}_k \odot \tilde{\mathbf{c}}_k$ with respect to $\mathbf{c}_{k-1}$ yields $\frac{\partial \mathbf{c}_k}{\partial \mathbf{c}_{k-1}} = \text{diag}(\mathbf{f}_k)$. Applying the chain rule across $T - t$ steps gives the product of diagonal matrices $\prod_{k=t+1}^T \text{diag}(\mathbf{f}_k)$. Setting $\mathbf{f}_k = \mathbf{1}$ reduces the product to the identity matrix $I$.
\end{proof}

\section{Complexity Analysis}

\begin{itemize}
    \item \textbf{Time Complexity:} $\mathcal{O}(4 \cdot (d_h d_x + d_h^2))$ FLOPs per time step (approximately $4\times$ the parameter and computational cost of a vanilla RNN).
    \item \textbf{Space Complexity:} $\mathcal{O}(2 \cdot d_h)$ memory to store state vectors $(\mathbf{h}_t, \mathbf{c}_t)$.
\end{itemize}

\section{Exactness and Error Analysis}

The forget gate bias initialization trick ($\mathbf{b}_f \gets 1.0$) ensures $\sigma(\mathbf{b}_f) \approx 0.73 > 0.5$ at initialization, guaranteeing memory persistence by default during early training phases.

\section{Worked Numerical Example}

Let $d_x = 1, d_h = 1, x_1 = 1.0, h_0 = 0.0, c_0 = 0.0$.
Parameters:
$W_f = [0.0], U_f = [0.0], b_f = [2.0] \implies \text{raw}_f = 2.0 \implies f_1 = \sigma(2.0) = 0.8808$.
$W_i = [1.0], U_i = [0.0], b_i = [0.0] \implies \text{raw}_i = 1.0 \implies i_1 = \sigma(1.0) = 0.7311$.
$W_c = [1.0], U_c = [0.0], b_c = [0.0] \implies \text{raw}_c = 1.0 \implies \tilde{c}_1 = \tanh(1.0) = 0.7616$.
$W_o = [0.0], U_o = [0.0], b_o = [0.0] \implies \text{raw}_o = 0.0 \implies o_1 = \sigma(0.0) = 0.5000$.

\textbf{State Updates:}
\begin{itemize}
    \item $c_1 = 0.8808 \times 0.0 + 0.7311 \times 0.7616 = 0.0 + 0.5568 = 0.5568$.
    \item $\tanh(c_1) = \tanh(0.5568) \approx 0.5055$.
    \item $h_1 = 0.5000 \times 0.5055 = 0.25275$.
\end{itemize}

\section{Numerical Considerations and Edge Cases}

\begin{itemize}
    \item \textbf{Fused Kernel Implementation:} Concatenating weights into a single $[4d_h, d_x + d_h]$ matrix enables high-throughput GEMM execution on modern GPUs.
    \item \textbf{Cell Value Explosion:} Without output tanh bounding ($\tanh(\mathbf{c}_t)$), cell states $\mathbf{c}_t$ can grow linearly without bound, eventually causing numerical overflow.
\end{itemize}

\section{References}
\begin{enumerate}
    \item Hochreiter, S., & Schmidhuber, J. (1997). Long short-term memory. \emph{Neural Computation}, 9(8), 1735--1780.
    \item Gers, F. A., Schmidhuber, J., & Cummins, F. (2000). Learning to forget: Continual prediction with LSTM. \emph{Neural Computation}, 12(10), 2451--2471.
\end{enumerate}

\end{document}
